\begin{table}[H]\centering
\caption{Leituras do teste isolado de M1 (cm)}\label{tab:m1-leituras}
\small\begin{tabular}{ccrrrrrr}\toprule
Comando & Ciclo & $X_i$ & $Y_i$ & $X_f$ & $Y_f$ & $X_r$ & $Y_r$\\\midrule
+65 & 1 & 0,30 & 44,50 & 3,10 & 44,30 & 1,00 & 43,80\\
+65 & 2 & 1,00 & 43,80 & 3,20 & 44,10 & 1,10 & 44,10\\
+65 & 3 & 1,10 & 44,10 & 3,00 & 44,50 & 1,30 & 44,30\\
-65 & 1 & -0,10 & 44,40 & -1,90 & 44,10 & -0,10 & 44,30\\
-65 & 2 & -0,10 & 44,30 & -2,00 & 43,90 & -0,20 & 44,10\\
-65 & 3 & -0,20 & 44,10 & -1,80 & 44,20 & -0,10 & 43,90\\
\bottomrule\end{tabular}\par\smallskip
{\footnotesize Fonte: Registros do autor. $i$: início da ida; $f$: final; $r$: retorno. Nos ciclos 2 e 3, o início corresponde ao retorno anterior.}\end{table}
\begin{table}[H]\centering
\caption{Deslocamento da ponta por ida no teste isolado de M1 (cm)}\label{tab:m1-deslocamento}
\small\begin{tabular}{crrr}\toprule Comando & Média $\Delta X$ & Média $\Delta Y$ & Média $\|\Delta\boldsymbol p\|$\\\midrule
+65 & 2,30 & 0,17 & 2,32\\
-65 & -1,77 & -0,20 & 1,79\\
\bottomrule\end{tabular}\par\smallskip
{\footnotesize Fonte: Elaborado pelo autor. A distância é a norma entre início e fim de cada ida, não o comprimento da trajetória.}\end{table}
